\section{Related Work}
\subsection{English-Pivot Reasoning}
대규모 언어 모델은 Chain-of-Thought (CoT) 프롬프팅~\citep{cot-neurips2022}을 통해 강력한 추론 능력을 보여주었으나, 이러한 능력은 언어에 따라 상당한 편차를 보인다. 다수의 선행 연구들은 비영어권 언어에서 추론 성능이 지속적으로 저하됨을 보고하였으며, 이러한 격차는 특히 수학적·논리적 추론 과제에서 두드러지게 나타난다~\citep{shi2022language, mega-emnlp2023, wang2025polymath}.

이러한 성능 차이를 완화하기 위해 영어를 피벗(pivot) 언어로 활용하는 프롬프팅 전략이 제안되었다. CLP~\citep{clp-emnlp2023}와 XLT~\citep{xlt-emnlp2023}는 비영어 질문을 영어로 변환하여 의미를 파악하게 하거나 추론 과정을 영어로 수행함으로써 성능 향상을 입증하였다. 후속 연구들 또한 비영어 입력을 영어로 번역하여 처리하는 것이 다양한 이해 및 추론 작업에 효과적임을 보여주었다~\citep{trans-allneed-2025naacl}.

그러나 이러한 접근법은 최종 답변이 영어로 생성될 경우 비영어권 사용자가 결과를 직관적으로 이해하거나 검증하기 어려워, 실용적인 측면에서 interpretability이 떨어진다는 한계를 가진다~\citep{multilingual-prompting-emnlp2025}.

\subsection{Structured Reasoning and Language Configuration} 
단순한 언어 변환을 넘어, 추론 과정을 명시적으로 구조화하여 복잡한 문제를 해결하려는 연구들이 제안되어 왔다. Skeleton-of-Thought~\citep{ning2023skeleton}와 Plan-and-Solve~\citep{wang2023plan}는 문제 해결을 위한 뼈대나 계획을 사전에 제시함으로써 추론의 일관성을 향상시킬 수 있음을 보였다. 이러한 접근을 바탕으로 다양한 구조 인식 프롬프팅으로 확장되었다~\citep{zhou2022least,khot2022decomposed,yao2023tree}.

다국어 환경에서도 유사한 시도가 이루어졌으나, 대부분 구조화된 정보 제공, 추론, 최종 출력까지의 언어 구성을 주로 영어로 고정하여 실험하였다. 예를 들어, \citet{qi2025sot}는 영어로 변수와 지식을 구조화하는 방식을 제안하였고, \citet{li2024eliciting}는 코드 형태의 표현을 활용하여 단계적 추론을 유도하였다. 그러나 스켈레톤이 표현되는 언어와 실제 추론 및 답변 언어의 조합이 성능에 미치는 영향은 체계적으로 분석되지 않았다.

본 연구는 이러한 구조적 추론 가이드를 \emph{스켈레톤}이라 통칭하며, 실험적 통제를 \citet{qi2025sot}의 구조를 베이스라인으로 채택한다. 우리는 스켈레톤이 표현되는 `언어'를 조작 변수로 설정하여, 다양한 언어 조합이 다국어 추론 성능에 미치는 영향을 심층 분석한다. 이를 통해 본 연구는 기존 문헌에서 다루어지지 않았던 언어 구성을 다국어 추론 시스템의 핵심 설계 변수로 정식화한다.



% 본 연구에서는 이러한 구조화된 추론 프롬프팅에서 사용되는 명시적 추론 구조를 스켈레톤(skeleton)이라는 상위 개념으로 통칭한다. 중요한 점은, 본 논문이 새로운 스켈레톤이나 프롬프팅 방식을 제안하는 것이 아니라, 구조와 형식은 고정한 채 스켈레톤이 표현되는 언어가 다국어 추론 성능에 미치는 영향을 분석한다는 것이다. 이를 통해 기존 연구에서 충분히 다루어지지 않았던 언어 구성(language configuration)을 핵심 설계 변수로 정식화한다.


% 단순한 언어 변환을 넘어, 추론 과정을 명시적으로 구조화하여 복잡한 문제를 해결하려는 시도 또한 활발히 이루어졌다. Skeleton-of-Thought~\citep{ning2023skeleton}나 Plan-and-Solve~\citep{wang2023plan}와 같은 연구는 문제 해결을 위한 뼈대(skeleton)나 계획을 먼저 생성하는 것이 모델의 논리적 일관성을 높이고 오류를 줄이는 데 효과적임을 보였다. 다국어 환경에서도 \citet{qi2025sot}는 영어로 변수와 지식을 구조화하는 방식을 제안하였고, \citet{li2024eliciting}는 코드(Code) 형식을 빌려 단계적 추론을 유도하였다.

% 선행 연구들은 Target Language로 직접 추론(CoT)을 수행할 때 발생하는 개념적 적용 오류(conceptual application error)~\citep{long-cot-2025arxiv}를 지적하면서도, 
% 하지만 기존의 skeleton과 같은 구조화된 방법론들은 정보를 모두 영어로 제공하거나, 입력부터 출력까지의 언어 구성을 단일하게 고정하여 실험하였다는 한계가 있다. 이를 완화하기 위한 수단으로서 스켈레톤 언어와 최종 답변 언어의 조합이 미치는 영향을 체계적으로 분석하지 않았다.


% 본 연구는 이러한 공백을 다루며, 영어 스켈레톤의 논리적 강점을 활용하면서도 Target Language 답변의 효용성을 유지하기 위한 최적의 언어 구성을 탐구한다는 점에서 기존 연구들과 차별화된다.